\documentclass{resume} % Use the custom resume.cls style

\usepackage[left=0.4 in,top=0.4in,right=0.4 in,bottom=0.4in]{geometry} % Document margins
\newcommand{\tab}[1]{\hspace{0.0\textwidth}\rlap{#1}} 
\newcommand{\itab}[1]{\hspace{0em}\rlap{#1}}
\name{Utkarsh Gupta} % Your name
% You can merge both of these into a single line, if you do not have a website.
\address{+91 9731802567 \\ Bengaluru, India \\ \href{mailto:utkarsh348@gmail.com}{utkarsh348@gmail.com} \\ \href{https://www.linkedin.com/in/utkarsh-gupta-0a621a191/}{linkedin.com/utkarsh-gupta}} 

\begin{document}

%----------------------------------------------------------------------------------------
%	OBJECTIVE
%----------------------------------------------------------------------------------------

%----------------------------------------------------------------------------------------
%	EDUCATION SECTION
%----------------------------------------------------------------------------------------



%----------------------------------------------------------------------------------------
% TECHINICAL STRENGTHS	
%----------------------------------------------------------------------------------------
\begin{rSection}{SKILLS}

\begin{tabular}{ @{} >{\bfseries}l @{\hspace{6ex}} l }
Languages &    Kotlin, Java, Python, TypeScript, Golang, SQL, C \\
Technologies & Kubernetes, Docker, Apache Spark, Apache Kafka, Apache Hadoop, Linux/UNIX \\
Data analysis Skills &  Exploratory Data Analysis, Supervised Learning, NLP, Computer Vision \\
Database and Cloud & MongoDB, DynamoDB, PostgreSQl - (SQL and NoSQL), AWS \\
Observability & Prometheus, Splunk \\
Frameworks/Libraries & Tensorflow, Keras, OpenCV, SpringBoot, Spring, Pandas \\
Back End Skills &  Microservice Architecture, REST API development, Secure Coding\\
LLM and Agentic AI & RAG, MCP Servers, MCP Clients, LLM SDKs \\
\end{tabular}
\end{rSection}

\begin{rSection}{EXPERIENCE}
 \itemsep-1em
\textbf{Software Engineer 2 - Intuit} \vspace{-5pt} \hfill Aug 2023 - Present
 \begin{itemize}
 \itemsep-0.5em
     \item Developed the OAuth2 authentication and authorization platform for Identity using SpringBoot and Kotlin with higher reliability at 99.999\%. 
     \item Designed token minting APIs for better fault tolerance and resiliency resulting in 65\% fewer false failures.
     \item Led the development of SAML protocols, saving more than \$400,000 in third-party contracts. 
     \item Led the development of BYOI(Bring Your Own Identity) SCIM Features enabling \$5 million in revenue.
     \item Developed AI enabled co-pilots, leveraging Model Context Protocol servers and RAG systems to boost developer productivity by 40\% and reducing MTTR by 30\% 
 \end{itemize}
\textbf{Co-op Intern - Intuit} \vspace{-6pt} \hfill Jan 2023 - Jun 2023
 \begin{itemize}
  \itemsep-0.5em
     \item Developed the system for seamless identity and access management (SCIM) for critical security requirements, using SpringBoot and Kotlin that delivered a 35\% faster onboarding experience. 
 \end{itemize}
 \textbf{Software Engineer Intern - Mosaic Wellness} \vspace{-6pt} \hfill Jun 2022 - July 2022
 \begin{itemize}
  \itemsep-0.5em
     \item Enhanced the backend system for appointment bookings and reports using TypeScript for 3x improvement in throughput and 50\% reduction in latency.
     \item Optimised SQL queries for memory efficiency and refactored queries to protect against injection attacks.
 \end{itemize}
\end{rSection} 
\begin{rSection}{Education}

{\bf Bachelor of Technology in Computer Science and Engineering}, PES University \hfill {2019 - 2023} \\
{8.9 GPA, Teaching Assistant for UE20CS333 Image Processing and Computer Vision
}
%Minor in Linguistics \smallskip \\
%Member of Eta Kappa Nu \\
%Member of Upsilon Pi Epsilon \\


\end{rSection}

%----------------------------------------------------------------------------------------
%	WORK EXPERIENCE SECTION
%----------------------------------------------------------------------------------------

\begin{rSection}{PROJECTS}
\vspace{-1em}
 \itemsep-0.5em
\item \textbf{Automatic Film Soundtrack generation} \vspace{-5pt}
\begin{itemize}
 \itemsep-0.7em
    \item Developed Autoencoders and Transformers to generate multi-dimensional sentiment-conditioned music. 
    \item Performed Sentiment analysis on the Valence-Arousal scale for scripts and successfully achieved transfer of attribute characteristics using Variational Autoencoders. 
    \item Published our work in the Generative AI Workshop at  \href{https://dl.acm.org/doi/10.1145/3639856.3639891}{AI-ML Systems 2023} and \href{https://arxiv.org/abs/2401.07084}{Machine Learning for Audio Workshop, NeurIPS 
2023}
\end{itemize}

\item \textbf{UK traffic data accidents' hotspot detection and analysis} \vspace{-5pt}
\begin{itemize}
 \itemsep-0.7em
 \item Used Kernel density estimation (KDE) plots to analyze hotspots of accident-prone areas weighed against severity over the years to understand the evolution of these dangerous zones. 
 \item Supervised learning algorithms were used to predict the accident severity. Later published at: \href{https://ieeexplore.ieee.org/document/10037449}{IBSSC 2022}.
 \end{itemize}
\item \textbf{Distributed File System Modeled On Hadoop}\vspace{-5pt} 
\begin{itemize}
 \itemsep-0.7em
 \item An implementation of a multithreaded DFS with namenodes and datanodes leveraging WebSockets for communication. The system is also capable of performing map-reduce tasks, caching chunks and is fault-tolerant to namenode and datanode failures.
   \end{itemize}

\item \textbf{Streaming Twitter via Spark and Kafka} \vspace{-5pt} 
\begin{itemize}
 \itemsep-0.7em
 \item Streaming Twitter data to Spark and publishing aggregated data through Kafka. Dynamically creating topics based on content of tweets.
  \end{itemize}

\end{rSection} 

%---------------------------------------------------


\end{document}
